% ====================================================================
\begin{abstract}
\noindent
I measured the electron spin resonance (ESR) of DPPH and \ce{MnCl2} with an X-band reflection spectrometer
(\qty{9.36}{GHz}). At the sample position the magnet's Hall probe reads too low, by \qty{8.5}{\percent} on the
analog output and \qty{2.3}{\percent} on the display, so I calibrated the field scale with the precisely known $g$-factor
of DPPH. The DPPH line has the peak-to-peak width $\Bpp=\qty{0.240(15)}{mT}$, corresponding to a spin-spin
relaxation time $T_2=\qty{27(2)}{ns}$. The signal grows linearly with the microwave power
($S_\mathrm{pp}\propto P^{\num{0.98(3)}}$), consistent with a square-law diode and without signs of saturation.
Up to \qty{1}{V} modulation voltage the signal grows proportionally. Since the line does not broaden even where the
amplifier limits (effectively \qty{1.8}{V}), the modulation amplitude at \qty{1}{V} is $\Bm<\qty{0.059}{mT}$.
Moving the plunger in the reference arm by \qty{5.60}{mm} turns absorption into dispersion. The phase shift of
\qty{-88(13)}{\degree} gives the guide wavelength $\lamg=46^{+8}_{-6}\,\unit{mm}$ and the cutoff frequency
$\nu_\mathrm c=6.7^{+0.8}_{-1.1}\,\unit{GHz}$, consistent with the WR-90 waveguide (\qty{44.9}{mm}, \qty{6.56}{GHz}).
\ce{MnCl2} shows a single Lorentzian line at $g=\num{2.058(11)}$ with the half width at half maximum (HWHM)
$\Bhalf=\qty{88.8(23)}{mT}$ ($T_2=\qty{62(2)}{ps}$). This fits fast exchange between the \ce{Mn^2+} ions, which averages out the dipolar
fields of the neighbors and the hyperfine structure. The copper sample, measured with a detuned source, gives no result.
\end{abstract}
